\documentclass[11pt]{article}
\usepackage{icmlko}
\begin{document}

\title{열여덟 칸 --- 판 꼴찌 도메인은 덜 배운 게 아니라 자가 굵었다}
\author{월드모델 실험실}
\date{노트 454 · 2026년 8월 2일}
\maketitle

\begin{abstract}
같은 자리를 세 노트가 연달아 가리켰다. 노트 444 는 CN 을 풀면서
\emph{``모형이 아니라 태거''}로 끝냈고, 노트 445 는 시장팝업의 같은 시기
천장이 $-0.0901$ 이라 \emph{``가진 게 전부 빌린 것''}이라 적었고, 노트 453
은 그 도메인만 의존도가 $0.0067$ 로 유독 낮은 것을 보고 \emph{``흔들
신호가 없다''}고 했다. 셋 다 관찰이었다.

이번엔 아주 싼 것을 쟀다. \textbf{축이 행을 몇 칸으로 가르나.} 행마다
축 값과 표시자를 묶어 칸을 세고, \textbf{칸 평균을 오라클로 주었을 때의
스피어만}을 그 축들이 낼 수 있는 상한으로 읽는다.

\begin{center}
\begin{tabular}{lrrrr}
\hline
도메인 & 행 & 칸 & 행/칸 & 오라클 상한 \\
\hline
아홉 도메인 & --- & $\approx$행 & $1.0$ & $\approx +1.00$ \\
만화 & $5{,}905$ & $3{,}136$ & $1.9$ & $+0.8374$ \\
팝업 & $89$ & $70$ & $1.3$ & $+0.8250$ \\
\textbf{시장팝업} & $249$ & $\mathbf{18}$ & $\mathbf{13.8}$ & $\mathbf{+0.2505}$ \\
\hline
\end{tabular}
\end{center}

\textbf{시장팝업 $249$행이 $18$칸에 들어간다. 그중 한 칸에 $162$행
($65\%$)이 있다.} 그 칸 안의 라벨 SD 가 $0.544$ 로 도메인 전체 폭과
같다 --- \emph{축이 삼분의 이에 대해 아무 말도 안 한다.}

$-0.0901$ 은 모형이 못해서 나온 게 아니다. \textbf{상한이 $+0.2505$ 인데
$18$칸에 $249$행으로 교차검증을 하니 과적합이 그 아래로 끌어내린 것}이다.
그리고 판 점수 $+0.1429$ 는 그 상한의 \textbf{$57\%$} --- \emph{이
도메인은 덜 배운 게 아니다.}
\end{abstract}

\begin{figure}[t]
\centering
\includegraphics[width=\textwidth]{figs/f1.pdf}
\caption{\textbf{가} --- 행/칸. 열한 도메인 중 축이 행을 못 가르는 것은
하나뿐이다. \textbf{나} --- 노트 453 이 못 푼 물음의 답. 의존도는
행 수가 아니라 칸 수를 따라간다.}
\end{figure}

\section{노트 453 의 열린 물음이 닫힌다}

노트 453 은 의존도(남만 흔들 때 내 점수의 SD)를 재고 \emph{행 수로는
설명이 안 된다}며 끝냈다 --- 아이돌 $122$행과 시장팝업 $123$행이 거의
같은 크기인데 의존도가 $6.7$배 달랐고, 행 수와의 스피어만은 $-0.200$
이었다.

\textbf{칸 수로는 설명된다.}

\begin{center}
\begin{tabular}{lrrr}
\hline
고정 도메인 & 학습 행 & 칸 & 의존도 SD \\
\hline
시장팝업 & $123$ & $\mathbf{18}$ & $0.0067$ \\
팝업 & $24$ & $70$ & $0.0349$ \\
아이돌 & $122$ & $173$ & $0.0447$ \\
도서 & $80$ & $236$ & $0.0398$ \\
\hline
\end{tabular}
\end{center}

\textbf{행 수 $-0.200$ 대 칸 수 $+0.800$.} 아이돌과 시장팝업은 행이
$122$ 대 $123$ 으로 같은데 \textbf{칸이 $9.6$배, 의존도가 $6.7$배}다.

기계도 말이 된다. 빌려주는 쪽을 흔들면 모델이 각 칸에 매기는 값이
움직인다. 칸이 $18$개뿐이면 움직일 자리가 $18$개고, 그중 하나에 행의
$65\%$ 가 몰려 있으니 순위가 거의 안 바뀐다. \textbf{칸은 남이 나를
흔들 수 있는 손잡이의 개수다.}

\section{그래서 무엇을 해야 하나}

\textbf{시장팝업에는 모형으로 할 일이 거의 없다.} 상한이 $+0.2505$ 이고
이미 $57\%$ 를 받고 있으니, 모형을 아무리 잘 만들어도 판에서 벌 수 있는
것은 $126/3{,}430 \times (0.25 - 0.14) = +0.004$ 가 최대다. 나머지는
\textbf{그 도메인에서 값이 갈리는 축을 만들어 주는 일} --- 태거 갈래다.

이것이 지금 열려 있는 갈래 중 \emph{``과제 78 시장 레코드 20건 사람
태깅''}이 왜 그 자리에 있는지 설명한다. 사람이 스무 건을 태깅하면
확인할 것은 정확도가 아니라 \textbf{그 스무 건이 몇 칸으로 갈리는가}다.

\section*{규약 --- 칸 수를 모든 실행에 찍는다}

축을 더할 때 우리는 늘 \emph{덮음}(관측 비율)을 봤다. 시장팝업의 세 축은
덮음이 $100\% \cdot 100\% \cdot 93\%$ 로 흠잡을 데가 없다. \textbf{덮음이
가득 차 있어도 값이 두 가지뿐이면 아무것도 못 가른다.}

\texttt{lab/loop.py} 의 \texttt{\_cell\_census} 로 모든 실행에 자동으로
찍는다 --- 도메인마다 칸 수 · 행/칸 · 오라클 상한, 그리고 \textbf{상한이
$0.9$ 아래인 도메인} 목록. 거부권은 아니고 이름표다(노트 415 · 401 과
같은 꼴).

\section*{남은 것}

\begin{itemize}
\item \textbf{진짜 덜 받는 도메인은 따로 있다.} 상한이 $\approx 1$ 인
도메인들에서는 이 비가 판 점수 자체라 뜻이 없지만, 적어도 시장팝업이
\emph{제일 덜 받는 도메인이 아니라는 것}은 확실하다. 펀딩 $+0.3689$ 와
도서 $+0.3746$ 은 상한이 $1.0$ 인데도 낮다.
\item \textbf{칸을 늘리면 정말 오르나.} 시장팝업에 값이 갈리는 축을
하나만 넣어도 칸이 몇 배가 된다 --- 사람 태깅 스무 건으로 먼저
\emph{갈릴 수 있는지}만 확인할 수 있다.
\item \textbf{일곱 빌려주는 도메인의 순수 의존도} --- 이제 칸 수로
예보가 되니, 재기 전에 예측을 적어 두고 잴 수 있다.
\end{itemize}

\end{document}
